% ===================================================================
%  TABLEAU N° 16 — Le Grand Magasin. — Le Bazar. — Les Jouets.
%  Feuillets 88 a 92 du fac-simile = folios 86 a 90.
%  Correspondance : folio imprime = numero de feuillet - 2.
%
%  Ceci est le dernier tableau du livret. Le feuillet 93 (« TABLE DES
%  MATIERES ») n'en fait pas partie et n'est pas releve ici.
%
%  Les cles « %%K » sont les MEMES que dans texto/io/25-tabelo-16.tex.
%  Le francais ne connait aucun des six sous-titres que l'ido
%  intercale (« Bela vidaji. », « Soldati e milito. », « Donacaji por
%  novyaro. », « Gapado cirkum butiko. », « Donacajo de mea avulo. ») :
%  Rochelle laisse courir la numerotation 1 a 18 sans intertitre, la
%  ou Guignon a ajoute ces repères. Ce n'est pas une lacune du releve.
%
%  ATTENTION — LE FAC-SIMILE FRANCAIS EST UNE PRISE DE VUE, contraste
%  inegal d'une page a l'autre.
%
%  Coquille conservee (§7 de la consigne) : au feuillet 88 (folio 86),
%  alinea 3, l'imprime porte « un singe (10 » en fin de ligne, SANS la
%  parenthese fermante — verifie a l'agrandissement (2600 px) : la
%  marge blanche suit immediatement le « 0 », aucune parenthese n'a
%  ete imprimee. Le numero de legende (10) est confirme par
%  texto/io/25-tabelo-16.tex (« simio (10) »).
%
%  Au feuillet 92 (folio 90), les renvois lettres (a) a (f) du tableau
%  d'astronomie sont composes en ITALIQUE dans le fac-simile francais
%  (contrairement aux renvois chiffres, droits) : klavize \textit{}.
%
%  La note en bas du feuillet 91 (folio 89), « (1) Voir Tableau 6. »,
%  est composee en gras et NON en italique (a la difference de la
%  note correspondante du livret ido, en italique) ; le chiffre « 6 »,
%  peu net a l'agrandissement, est confirme par le contenu meme du
%  Tableau N° 6 (« La Maison intérieure. — Les Meubles. ») et par la
%  note homologue de texto/io/25-tabelo-16.tex (« Videz la tabelo
%  N° 6. »). Le renvoi « (1) » dans le corps du texte est lui-meme
%  compose en taille normale, non en exposant, a la difference des
%  renvois numerotes de legende.
% ===================================================================

% --- Feuillet 88 (folio 86, ouverture de tableau) -------------------
\begin{VUpage}[88]{86}
%%K t16-apar-1 apar
\VUsaut{6.0mm}
\VUtitre{15.0pt}{60}{TABLEAU N\textsuperscript{o} 16}
\VUsaut{1.4mm}
\VUtitre{11.4pt}{0}{Le Grand Magasin. --- Le Bazar. --- Les Jouets.}
\VUsaut{2.2mm}

%%K t16-01-1 p
\VUblancAlinea
1. --- Le nouvel an approche. Les grands magasins\nl
ont préparé leurs étalages de \VUgras{jouets}\textsuperscript{(1)} et de \VUgras{cadeaux}\nl
pour les \VUgras{étrennes.}

%%K t16-01-2 p
\VUblancAlinea
J'ai demandé à mon grand-père de me conduire au\nl
bazar pour voir cette belle exposition et il y a consenti.

%%K t16-02-1 p
\VUblancAlinea
2. --- Nous nous sommes d'abord arrêtés devant de\nl
jolies panoplies, qui comprennent un \VUgras{fusil}, un \VUgras{képi,}\nl
un \VUgras{sabre}, un \VUgras{ceinturon} et une paire d'\VUgras{épaulettes,}\nl
ou bien un \VUgras{casque}, une \VUgras{cuirasse}\textsuperscript{(3)} et un \VUgras{pis}\cc
\VUgras{tolet}\textsuperscript{(4)}. Tout cela m'intéressait beaucoup plus que\nl
les \VUgras{lanternes magiques}\textsuperscript{(2)}.

%%K t16-03-1 p
\VUblancAlinea
3. --- Il y avait dans le même rayon un \VUgras{faction}\cc
\VUgras{naire}\textsuperscript{(5)} dans sa \VUgras{guérite} ; il avait l'air de monter\nl
la garde devant toute une ménagerie en carton. Que\nl
d'animaux il y avait là ! Un \VUgras{perroquet}\textsuperscript{(6)} sur son\nl
perchoir, un \VUgras{rhinocéros}\textsuperscript{(7)} avec une corne sur le\nl
nez, un \VUgras{renne}\textsuperscript{(8)}, une \VUgras{autruche}\textsuperscript{(9)}, un \VUgras{singe}\textsuperscript{(10}\nl
croquant une noix, un \VUgras{renard}\textsuperscript{(11)} qui emportait une\nl
poule, un \VUgras{crocodile}\textsuperscript{(12)} qui ouvrait ses énormes mâ\cc
choires, un \VUgras{chameau}\textsuperscript{(13)}, un élégant petit \VUgras{lévrier}\textsuperscript{(14)},\nl
une \VUgras{panthère}\textsuperscript{(15)}, puis deux bêtes que je ne connais\cc
sais pas et qui sont, paraît-il, une \VUgras{belette}\textsuperscript{(16)} et une\nl
\VUgras{hyène}\textsuperscript{(17)}. Il y avait encore un \VUgras{léopard}\textsuperscript{(18)} et un\nl
\VUgras{loup}\textsuperscript{(21)} ; tout auprès se trouvaient de mignons petits\nl
\VUgras{rats}\textsuperscript{(19)} et de minuscules petites \VUgras{souris}\textsuperscript{(20)} en caout\cc
chouc. Enfin un \VUgras{hippopotame}\textsuperscript{(22)} et un \VUgras{droma}\cc
\VUgras{daire}\textsuperscript{(23)} bossu terminaient la collection. Je serais\nl
resté là pendant des heures s'il n'y avait eu à côté un\nl
magnifique camp, rempli de \VUgras{soldats de plomb}\textsuperscript{(29)},\nl
qui attira mon attention.

%%K t16-04-1 p
\VUblancAlinea
4. --- Mon grand-père, qui est \VUgras{invalide}\textsuperscript{(24)} et qui\nl
a dû quitter l'armée à la suite d'une \VUgras{blessure} qui
\parplein
\end{VUpage}

% --- Feuillet 89 (folio 87) -----------------------------------------
\begin{VUpage}[89]{87}
%%K t16-04-1 p suite
\VUcontinue
lui valut la \VUgras{médaille militaire}, aime beaucoup à me\nl
raconter les \VUgras{campagnes} auxquelles il a pris part.\nl
Il fut heureux de m'expliquer les différents mouve\cc
ments de la petite armée en miniature. Un groupe de\nl
vrais soldats qui visitaient le bazar, un \VUgras{soldat du}\nl
\VUgras{génie}\textsuperscript{(25)}, un \VUgras{chasseur alpin}\textsuperscript{(26)} coiffé d'un \VUgras{béret,}\nl
un \VUgras{sergent-major}\textsuperscript{(27)} d'infanterie et un \VUgras{maréchal}\nl
\VUgras{des logis}\textsuperscript{(28)} d'artillerie, qui avait un \VUgras{galon} d'or sur\nl
sa manche, s'arrêtèrent près de nous pour écouter les\nl
explications.

%%K t16-05-1 p
\VUblancAlinea
5. --- Grand-père, tout fier d'avoir un si brillant\nl
auditoire, improvisa tout un plan de bataille.

%%K t16-05-2 p
\VUblancAlinea
La ville forte, avec ses \VUgras{remparts} et ses \VUgras{fossés,}\nl
était assiégée par l'\VUgras{armée ennemie} qui, après un\nl
long \VUgras{bombardement,} se disposait à donner l'\VUgras{assaut.}\nl
Mais les \VUgras{assiégés,} poussés à bout par la famine,\nl
avaient tenté une \VUgras{sortie} contre le \VUgras{camp} des \VUgras{assié}\cc
\VUgras{geants.} Après avoir repoussé l'\VUgras{attaque} dans un\nl
\VUgras{combat} acharné autour des \VUgras{tentes}, les assiégeants\nl
\VUgras{vainqueurs} lançaient leurs \VUgras{escadrons} à la \VUgras{pour}\cc
\VUgras{suite} des assaillants \VUgras{vaincus.} Ceux-ci battaient en\nl
\VUgras{retraite}, couvrant le \VUgras{champ de bataille} de \VUgras{morts}\nl
et de \VUgras{blessés.} Des \VUgras{brancardiers} emportaient ces\nl
derniers à l'\VUgras{ambulance} pour les faire soigner par\nl
le \VUgras{chirurgien.}

%%K t16-05-3 p
\VUblancAlinea
Malgré l'héroïque résistance de quelques \VUgras{batail}\cc
\VUgras{lons} qui s'étaient formés en \VUgras{carré}, la \VUgras{défaite} était\nl
complète, le reste de l'armée prenait la \VUgras{fuite}, aban\cc
donnant de nombreux \VUgras{prisonniers.} L'ennemi allait\nl
couronner sa \VUgras{victoire} en s'emparant de la ville. Ce\nl
serait peut-être la fin de la campagne et, après un\nl
\VUgras{armistice}, on pourrait signer le \VUgras{traité de paix.}

%%K t16-06-1 p
\VUblancAlinea
6. --- Les jeunes soldats, qui riaient en écoutant le\nl
pittoresque récit du vétéran, lui demandèrent quelques\nl
autres explications. Moi, je fis le tour du comptoir\nl
pour regarder une belle \VUgras{arbalète}\textsuperscript{(32)} et un \VUgras{arc}\textsuperscript{(33)}\nl
avec sa \VUgras{flèche}\textsuperscript{(31)}. Je trouvai là une autre collection\nl
d'animaux : deux \VUgras{oiseaux chanteurs}\textsuperscript{(34)}, un \VUgras{rossi}\cc
\parplein
\end{VUpage}

% --- Feuillet 90 (folio 88) -----------------------------------------
\begin{VUpage}[90]{88}
%%K t16-06-1 p suite
\VUcontinue
\VUgras{gnol} et une \VUgras{fauvette} mécaniques qui chantaient à\nl
plein gosier, et une série de bêtes bizarres, \VUgras{chauve}\cc
\VUgras{souris}\textsuperscript{(35)}, \VUgras{serpent boa}\textsuperscript{(36)}, \VUgras{tortue}\textsuperscript{(37)}, \VUgras{phoque}\textsuperscript{(38)},\nl
\VUgras{baleine}\textsuperscript{(39)} et \VUgras{requin}\textsuperscript{(40)} en baudruche.

%%K t16-07-1 p
\VUblancAlinea
7. --- Je ne m'arrêtai pas longtemps à les contem\cc
pler, car j'avais aperçu l'objet de mes rêves, un\nl
superbe \VUgras{théâtre}\textsuperscript{(41)} de marionnettes, avec de magni\cc
fiques \VUgras{décors} et un ravissant \VUgras{rideau} rouge. En\nl
\VUgras{scène}, deux \VUgras{acteurs} avaient l'air de jouer un \VUgras{rôle}\nl
dans une \VUgras{pièce} fort intéressante, car les petits \VUgras{spec}\cc
\VUgras{tateurs} de carton installés dans les \VUgras{loges} ou assis\nl
aux \VUgras{fauteuils} et au \VUgras{parterre} derrière le \VUgras{chef}\nl
\VUgras{d'orchestre,} applaudissaient vivement.

%%K t16-07-2 p
\VUblancAlinea
Malheureusement, ce théâtre coûtait très cher.

%%K t16-08-1 p
\VUblancAlinea
8. --- D'ailleurs, le choix des jouets était grand, car\nl
sur les vitrines toutes sortes de jouets s'entassaient :\nl
\VUgras{Arlequin}\textsuperscript{(42)}, \VUgras{Polichinelle}\textsuperscript{(43)}, \VUgras{Pierrot}\textsuperscript{(44)}, des\nl
\VUgras{hiboux}\textsuperscript{(45)} battant des ailes, des \VUgras{merles}\textsuperscript{(46)} siffleurs,\nl
des \VUgras{caniches} qui aboient, un \VUgras{phonographe}\textsuperscript{(47)} et\nl
des \VUgras{jeux de patience.} Un de ces jeux m'amusait\nl
beaucoup ; il représentait un \VUgras{combat naval}\textsuperscript{(48)} dans\nl
lequel un \VUgras{navire} est attaqué par des \VUgras{pirates} près\nl
d'un pays planté de \VUgras{cocotiers.}

%%K t16-09-1 p
\VUblancAlinea
9. --- Je pouvais contempler très facilement toutes\nl
ces merveilles, car il n'y avait que peu de monde dans\nl
le magasin, à peine quelques rares flâneurs, un\nl
\VUgras{dragon}\textsuperscript{(49)}, et un \VUgras{encaisseur}\textsuperscript{(50)} qui présentait une\nl
\VUgras{traite} à la \VUgras{caisse}\textsuperscript{(51)} principale.

%%K t16-10-1 p
\VUblancAlinea
10. --- Je demandai à mon grand-père à quoi ser\cc
vaient les gros livres placés sur des rayons derrière\nl
la \VUgras{caissière} ; il me répondit que c'étaient les livres\nl
de \VUgras{comptabilité.}

%%K t16-11-1 p
\VUblancAlinea
11. --- En quittant le rayon des jouets, nous som\cc
mes allés à la sellerie jeter un coup d'œil sur les\nl
\VUgras{selles}\textsuperscript{(53)}, les \VUgras{étriers}\textsuperscript{(54)}, les \VUgras{brides}\textsuperscript{(55)}, les\nl
\VUgras{mors}\textsuperscript{(56)} et les \VUgras{cravaches.} Pendant que le \VUgras{chef de}\nl
\VUgras{rayon}\textsuperscript{(57)} donnait quelques renseignements à grand
\parplein
\end{VUpage}

% --- Feuillet 91 (folio 89) -----------------------------------------
\begin{VUpage}[91]{89}
\VUnotes{143.2mm}{%
\VUgras{(1) Voir Tableau 6.}%
}

%%K t16-11-1 p suite
\VUcontinue
papa, j'ai été regarder l'étalage des \VUgras{porcelaines} et\nl
des \VUgras{cristaux}\textsuperscript{(58)}, où il y avait de jolis \VUgras{saladiers} et\nl
de fines \VUgras{théières}\textsuperscript{(59)}.

%%K t16-12-1 p
\VUblancAlinea
12. --- Pour aller au premier étage, nous passâmes\nl
derrière la vitrine des \VUgras{corsets}\textsuperscript{(60)}, près de la\nl
bonneterie, où étaient étalés de très beaux \VUgras{cale}\cc
\VUgras{çons}\textsuperscript{(63)}. A côté, une \VUgras{demoiselle de magasin}\textsuperscript{(61)}\nl
faisait choisir à une \VUgras{acheteuse}\textsuperscript{(62)} de jolis \VUgras{tissus}\textsuperscript{(64)}\nl
de soierie.

%%K t16-12-2 p
\VUblancAlinea
En montant l'escalier, j'ai remarqué que le magasin\nl
était décoré de \VUgras{lanternes}\textsuperscript{(65)} japonaises, de \VUgras{para}\cc
\VUgras{sols}\textsuperscript{(66)} et d'énormes \VUgras{palmiers}\textsuperscript{(70)}.

%%K t16-13-1 p
\VUblancAlinea
13. --- Au premier, il n'y avait pas grand'chose à\nl
voir : rien que des meubles (1), des \VUgras{coffres-forts}\textsuperscript{(67)},\nl
des \VUgras{commodes}\textsuperscript{(68)} et des \VUgras{voitures d'enfant}\textsuperscript{(69)}.\nl
Mais la vue d'ensemble du magasin était très amu\cc
sante.

%%K t16-14-1 p
\VUblancAlinea
14. --- A nos pieds, je voyais les instruments de\nl
musique : \VUgras{harpes}\textsuperscript{(71)}, \VUgras{guitares}\textsuperscript{(72)}, \VUgras{mandolines}\textsuperscript{(73)},\nl
\VUgras{cornets à piston}\textsuperscript{(74)}, \VUgras{clarinettes}\textsuperscript{(75)}, violons\nl
avec leur \VUgras{archet}\textsuperscript{(76)} et même une \VUgras{grosse caisse}\textsuperscript{(77)}.\nl
Un peu plus loin, au rayon de la papeterie, un \VUgras{étu}\cc
\VUgras{diant}\textsuperscript{(78)} marchandait une serviette au \VUgras{commis}\textsuperscript{(79)}.\nl
Près d'eux, je distinguai un bel \VUgras{alphabet}\textsuperscript{(80)}.

%%K t16-14-2 p
\VUblancAlinea
Au milieu du magasin se dressait un grand \VUgras{arbre}\nl
\VUgras{de Noël}\textsuperscript{(81)} tout garni de bougies, de bibelots et de\nl
jouets de toute sorte : \VUgras{chevaux de bois}\textsuperscript{(82)}, \VUgras{pou}\cc
\VUgras{pées}\textsuperscript{(83)}, etc.

%%K t16-14-3 p
\VUblancAlinea
La \VUgras{parfumerie}\textsuperscript{(84)}, avec ses \VUgras{eaux dentifrices}\nl
et ses \VUgras{parfums} variés, répandait une très bonne\nl
odeur qui montait jusqu'à nous.

%%K t16-15-1 p
\VUblancAlinea
15. --- Comme grand-père commençait à trouver le\nl
temps long, nous sommes redescendus par le grand\nl
escalier. Il s'arrêta un instant devant un tableau\nl
d'\VUgras{astronomie}\textsuperscript{(85)} représentant, m'a-t-il dit, des
\parplein
\end{VUpage}

% --- Feuillet 92 (folio 90) -----------------------------------------
\begin{VUpage}[92]{90}
%%K t16-15-1 p suite
\VUcontinue
\VUgras{comètes}\textsuperscript{\textit{(a)}}, des \VUgras{éclipses}\textsuperscript{\textit{(b)}} de lune et de soleil,\nl
une rose des vents avec les \VUgras{quatre points cardi}\cc
\VUgras{naux}\textsuperscript{\textit{(c)}} (le \VUgras{nord}, le \VUgras{sud}, l'\VUgras{est} et l'\VUgras{ouest}), une\nl
carte des deux \VUgras{hémisphères}\textsuperscript{\textit{(d)}}, les \VUgras{planètes}\textsuperscript{\textit{(e)}} et\nl
le \VUgras{système solaire}\textsuperscript{\textit{(f)}}. Je ne comprenais rien à tout\nl
cela, et je m'intéressais beaucoup plus à la livrée\nl
galonnée d'un \VUgras{garçon livreur}\textsuperscript{(86)} qui passait et aux\nl
jolis \VUgras{éventails}\textsuperscript{(87)} qui étaient près de moi, à côté des\nl
\VUgras{porte-monnaie}\textsuperscript{(88)} et des \VUgras{portefeuilles}\textsuperscript{(89)}.

%%K t16-16-1 p
\VUblancAlinea
16. --- J'admirai aussi sur le comptoir des \VUgras{plu}\cc
\VUgras{mes}\textsuperscript{(90)} et des \VUgras{boucles}\textsuperscript{(91)} de ceinture, les jolies\nl
\VUgras{fleurs artificielles}\textsuperscript{(94)} gracieusement disposées\nl
dans des corbeilles. Les \VUgras{violettes}, les \VUgras{giroflées,}\nl
les \VUgras{jacinthes}\textsuperscript{(95)}, les \VUgras{géraniums}\textsuperscript{(96)}, les \VUgras{lis}\textsuperscript{(97)} et\nl
les \VUgras{capucines}\textsuperscript{(98)} étaient très bien imités.

%%K t16-17-1 p
\VUblancAlinea
17. --- Je priai mon aïeul de m'acheter une des\nl
\VUgras{brosses à dents}\textsuperscript{(92)} qui se trouvaient dans une\nl
boîte près des \VUgras{épingles à cheveux}\textsuperscript{(93)}. Il y con\cc
sentit et j'allai moi-même payer mon emplette à la\nl
caisse, auprès de laquelle on préparait une \VUgras{expédi}\cc
\VUgras{tion}\textsuperscript{(100)} importante pour un \VUgras{client}\textsuperscript{(99)}.

%%K t16-18-1 p
\VUblancAlinea
18. --- Tout à coup une légère agitation se produisit\nl
au milieu d'un groupe de personnes arrêtées près des\nl
\VUgras{bronzes}\textsuperscript{(101)} artistiques. C'était un \VUgras{voleur}\textsuperscript{(102)} qui\nl
venait d'être pris en flagrant délit par un \VUgras{inspec}\cc
\VUgras{teur}\textsuperscript{(103)}, pendant qu'il essayait de dévaliser un\nl
monsieur. On fit appeler un agent de police pour\nl
procéder à son \VUgras{arrestation} et, au moment où nous\nl
sortions, le malfaiteur était emmené au poste.

\VUsaut{6.0mm}
\VUfilet{22mm}
\end{VUpage}
